The greatest eclipse is defined as the instant when the axis of the Moon's
shadow cone gets closest to the centre of the Earth in a solar eclipse. This
problem explores the geometry of this phenomenon, using the solar eclipse of
29\textsuperscript{th} May 1919 as an example, as it has great historical
significance for being the first time astronomers were able to
observationally verify general relativity. One of the scientific expeditions
to observe this eclipse took place in the Brazilian city of Sobral.

The two following tables show the Cartesian and spherical coordinates of the
Sun and the Moon at the time of the greatest eclipse. The system used for
these coordinates is right-handed and has the origin at the centre of the
Earth, the positive $x$-axis pointing towards the Greenwich meridian, and the
positive $z$-axis pointing towards the North Pole. For the rest of this
problem, this will be referred to as \textbf{System I}.

Spherical Coordinates:
\begin{center}
\begin{tabular}{c|c|c}
 & Centre of the Sun & Centre of the Moon \\
\hline
Radial Distance ($r$) & $1.516 \times 10^{11}$\,m
  & $3.589 \times 10^{8}$\,m \\
Polar Angle ($\theta$) & $68^\circ 29' 44.1''$ & $68^\circ 47' 41.6''$ \\
Azimuthal Angle ($\varphi$) & $-1^{\mathrm{h}}11^{\mathrm{m}}28.2^{\mathrm{s}}$
  & $-1^{\mathrm{h}}11^{\mathrm{m}}22.9^{\mathrm{s}}$ \\
\end{tabular}
\end{center}

Cartesian Coordinates:
\begin{center}
\begin{tabular}{c|c|c}
 & Centre of the Sun & Centre of the Moon \\
\hline
$x$ & $1.342 \times 10^{11}$\,m & $3.185 \times 10^{8}$\,m \\
$y$ & $-4.327 \times 10^{10}$\,m & $-1.025 \times 10^{8}$\,m \\
$z$ & $5.557 \times 10^{10}$\,m & $1.298 \times 10^{8}$\,m \\
\end{tabular}
\end{center}

For this problem, assume that the Earth is a perfect sphere.

\textbf{Note:} The spherical coordinates of a point $P$ are defined as
follows:
\begin{itemize}
\item Radial distance ($r$): distance between the origin ($O$) and $P$
  (range: $r \geq 0$).
\item Polar angle ($\theta$): angle between the positive $z$-axis and the
  line segment $OP$ (range: $0^\circ \leq \theta \leq 180^\circ$)
\item Azimuthal angle ($\varphi$): angle between the positive $x$-axis and
  the projection of the line segment $OP$ onto the $xy$-plane (range:
  $-12^{\mathrm{h}} \leq \varphi < 12^{\mathrm{h}}$)
\end{itemize}

\ioaafig{2024_TH_T10-f1.pdf}
% point P, radius r, polar angle theta and azimuthal angle phi (page 10 of
% the 2024 theory paper, top of page)

\textbf{Part I: Geographic Coordinates (25 points)}

\begin{description}
\item[(a)] (3 points) Determine the declination of the Sun and the Moon
  during the greatest eclipse for a geocentric observer.
\item[(b)] (3 points) Determine the right ascension of the Sun and the Moon
  at the time of the greatest eclipse for a geocentric observer. The local
  sidereal time at Greenwich at that same moment was
  $5^{\mathrm{h}}32^{\mathrm{m}}35.5^{\mathrm{s}}$.
\item[(c)] (4 points) Find a unit vector that indicates the direction of the
  axis of the Moon's shadow cone. This vector should point from the Moon to
  the vicinity of the centre of the Earth.
\item[(d)] (15 points) Determine the latitude and the longitude of the point
  where the axis of the Moon's shadow cone crosses the surface of the Earth
  during the greatest eclipse.
\end{description}

\textbf{Part II: Duration of the Totality (50 points)}

Precisely determining the duration of totality of a solar eclipse involves
complex calculations that would be beyond the scope of this problem. However,
it is possible to obtain a reasonable approximation for this value using the
two following assumptions:
\begin{itemize}
\item The size of the umbra on the surface of the Earth remains roughly
  constant throughout totality at a given location.
\item The velocity of the umbra on the surface of the Earth remains roughly
  constant throughout totality at a given location.
\end{itemize}

\begin{description}
\item[(e)] (10 points) Estimate the radius of the umbra during the greatest
  eclipse. In order to simplify the calculations, assume that the umbra is
  small enough that it can be considered approximately flat and that the axis
  of the Moon's shadow cone is extremely close to the centre of the Earth
  during the greatest eclipse.
\item[(f)] (3 points) Calculate the velocity of the Earth's rotation at the
  latitude of the centre of the umbra.
\item[(g)] (4 points) Determine the orbital velocity of the Moon at the
  instant of the greatest eclipse. Neglect the changes in the semi-major axis
  of the Moon's orbit.
\end{description}

For the remaining items of this problem, assume that the tangential velocity
of the Moon is roughly the same as the orbital velocity and neglect its
radial component.

In order to calculate the velocity of the umbra, it is convenient to define
two new additional right-handed coordinate systems. \textbf{System II} is
defined as follows:
\begin{itemize}
\item Origin ($O_{\mathrm{II}}$): position of the Moon at the instant of the
  greatest eclipse.
\item Positive $x$-axis: Tangent to the declination circle. Points
  eastwards.
\item Positive $y$-axis: Tangent to the meridian of right ascension. Points
  northwards.
\end{itemize}

\textbf{System III} is defined as follows:
\begin{itemize}
\item Origin ($O_{\mathrm{III}}$): centre of the umbra at the instant of the
  greatest eclipse.
\item Positive $x$-axis: Tangent to the latitude circle. Points eastwards.
\item Positive $y$-axis: Tangent to the meridian of longitude. Points
  northwards.
\end{itemize}

Note that in both systems, the $xy$-plane is tangent to the celestial sphere
at the position of the origin.

System III is similar to System II, with the only difference being that the
origin ($O_{\mathrm{III}}$) is at the centre of the umbra at the moment of
the greatest eclipse.

\begin{description}
\item[(h)] (14 points) Using System II, determine the velocity vector of the
  Moon during the greatest eclipse. Note that the intersection between the
  Celestial Equator and the lunar orbit that is closer to the position of the
  eclipse has a right ascension of
  $23^{\mathrm{h}}07^{\mathrm{m}}59.2^{\mathrm{s}}$.
\item[(i)] (10 points) Write the velocity vector of the Moon in System III.
  Note that in System I, the azimuthal angle difference between the positions
  of the origins $O_{\mathrm{II}}$ and $O_{\mathrm{III}}$ is negligible, so
  you should only take into account the difference in the polar angles.
\item[(j)] (6 points) Calculate the speed of the centre of the umbra along
  the surface of the Earth at the instant of the greatest eclipse.
\item[(k)] (3 points) Estimate the duration of the totality of the eclipse at
  the location with the coordinates found in part (d).
\end{description}
